\section{Определение оригинала и изображения по Лапласу. Теорема о
существовании изображения и его
аналитичности.}\label{ux43eux43fux440ux435ux434ux435ux43bux435ux43dux438ux435-ux43eux440ux438ux433ux438ux43dux430ux43bux430-ux438-ux438ux437ux43eux431ux440ux430ux436ux435ux43dux438ux44f-ux43fux43e-ux43bux430ux43fux43bux430ux441ux443.-ux442ux435ux43eux440ux435ux43cux430-ux43e-ux441ux443ux449ux435ux441ux442ux432ux43eux432ux430ux43dux438ux438-ux438ux437ux43eux431ux440ux430ux436ux435ux43dux438ux44f-ux438-ux435ux433ux43e-ux430ux43dux430ux43bux438ux442ux438ux447ux43dux43eux441ux442ux438.}

\textbf{Определение:} \emph{Оригиналом} называют комплекснозначную
функцию \(f(t)\) действительного переменного \(t\), если:

\begin{enumerate}
\def\labelenumi{\arabic{enumi}.}
\tightlist
\item
  \(f(t)=0\) при \(t<0\)
\item
  на каждом отрезке полуоси \(t\geq 0\) функция \(f(t)\) непрерывна,
  кроме, быть может, конечного числа точек разрыва первого рода;
\item
  Существуют такие действительные числа \(M>0\) и \(\alpha\), что
  \(\forall t\geq 0\) выполняется неравенство: \[
  |f(t)|\leq Me^{ \alpha t }
  \]
\end{enumerate}

\textbf{Определение:} \emph{Изображением} оригинала \(f(t)\) называют
комплекснозначную функцию \[
F(p)=\int_{0}^{+\infty}e^{ -pt }\cdot f(t)\,dt\quad(2)
\] комплексного переменного \(p\).

\textbf{Определение:} Интеграл (2) называют \emph{преобразованием
Лапласа} функции \(f(t)\).

Связь между оригиналом \(f(t)\) и его изображением \(F(p)\) записывают
так: \[
f(t)\risingdotseq F(p)\text{ или }F(p)\risingdotseq f(t)
\] \textbf{Теорема:} Пусть \(f(t)\) - оригинал. Тогда \(f(t)\) имеет
изображение \(F(p)=\int_{0}^{\infty}f(t)e^{ -pt }\,dt\) и \(F(p)\) -
аналитическая функция в области
\(G_{\alpha}=\{ p\in \mathbb{C}: \mathrm{Re}p>\alpha \}\).
